\documentclass[11pt]{article}

\usepackage[utf8]{inputenc}
\usepackage[T1]{fontenc}
\usepackage{lmodern}
\usepackage{geometry}
\usepackage{amsmath,amssymb,amsfonts,amsthm,mathtools}
\usepackage{bm}
\usepackage{enumitem}
\usepackage{microtype}
\usepackage{hyperref}
\usepackage{titlesec}
\usepackage{booktabs}
\usepackage{array}
\usepackage{setspace}
\usepackage{mathrsfs}
\usepackage{physics}

\geometry{margin=1in}
\hypersetup{colorlinks=true, linkcolor=black, urlcolor=blue, citecolor=black}
\setlist[enumerate]{topsep=0.4em,itemsep=0.35em}
\setlist[itemize]{topsep=0.3em,itemsep=0.25em}
\titleformat{\section}{\large\bfseries}{\thesection.}{0.5em}{}
\titleformat{\subsection}{\normalsize\bfseries}{\thesubsection.}{0.5em}{}
\setstretch{1.06}

\newcommand{\Aether}{\AE{}ther}
\newcommand{\Aflow}{\AE{}ther-flow}
\newcommand{\TheoryName}{\Aflow{} Interpretation of Relativity}
\newcommand{\TheoryProgram}{\Aether{} / \Aflow{} framework}
\newcommand{\Sub}{\mathcal{A}}
\newcommand{\BridgeMetric}{\mathfrak{g}}
\newcommand{\Lor}{\gamma_{\mathrm{L}}}

\newtheorem{definition}{Definition}
\newtheorem{proposition}{Proposition}
\newtheorem{remark}{Remark}
\newtheorem{corollary}{Corollary}
\newtheorem{theorem}{Theorem}

\title{\textbf{The \TheoryName{}}\\[0.4em]
\large Higher-Derivative Quartic Branch Compatibility with the Exact-Relativistic Recovery Line}
\author{}
\date{}

\begin{document}

\maketitle

\begin{abstract}
This manuscript performs the next fixed gate for the higher-derivative quartic branch after the positive flat-background exact-closure test and the bridge-metric matter-coupling gate. The scope is again narrow and explicit. The note does not yet promote the branch to a local curved-background continuation, and it does not claim a completed substrate derivation of Einsteinian gravity. It checks only whether the same minimal quartic branch remains compatible with the active exact-relativistic recovery line already fixed by the dynamics and relativistic-recovery manuscripts.

The result is positive at the present action-level and flat-branch quadratic scope, provided the surviving quartic metric self-energy is evaluated in the infrared regime for which it is parametrically subleading. On the decaying flat branch with
\[
m_S^2>0,
\qquad
\lambda_4>0,
\qquad
\widehat M_{IJ}\ \text{positive definite},
\]
the only surviving observer-accessible branch correction remains the metric scalar self-energy
\[
\Delta\mathcal L_{\mathrm{metric},4}^{(2)}
=
-\frac{1}{2}\Pi_\phi^{(4)}(\omega,\mathbf{k})\,|\phi(\omega,\mathbf{k})|^2,
\qquad
\Pi_\phi^{(4)}
=
m_S^2\sigma_0^2
\frac{\mathcal B_4(k^2)}
{m_S^2\omega^2+\mathcal B_4(k^2)},
\]
with
\[
\mathcal B_4(k^2)=\frac{\lambda_4}{M_*^2}k^4.
\]
Hence the reduced scalar operator satisfies
\[
\mathcal B_{4,0}=0,
\qquad
\mathcal B_{4,2}=0,
\]
so the GR weak-field/Newtonian sector is preserved through the benchmark \(k^0/k^2\) screen. Because the observer map, bridge metric, and universal matter route remain unchanged, the branch also preserves the single operational metric, the unique null cone, the proper-time rule, and the local special-relativistic limit already fixed in the active recovery manuscript. The remaining quartic correction is below observer-accessible relevance in the infrared regime
\[
k^2\ll M_*^2,
\qquad
\varepsilon_4(\omega,\mathbf{k})
\equiv
\frac{\lambda_4 k^4}{m_S^2M_*^2\omega^2}
\ll 1,
\]
for which
\[
\Pi_\phi^{(4)}
=
m_S^2\sigma_0^2\,\varepsilon_4
+ \mathcal O(\varepsilon_4^2).
\]
In that precise current-scope sense, the quartic branch is compatible with the active exact-relativistic recovery line. This is the clean scope that may be carried forward to a later local curved-background continuation; the ultrasoft static corner is not claimed here as already settled.
\end{abstract}

\section{Introduction}

The active exact-closure line of \TheoryName{} is already fixed at the observer-accessible level. The dynamics manuscript adopts a single Einsteinian metric with universal matter coupling. The relativistic-recovery manuscript then states explicitly that weak-field gravity, Newtonian reduction, null propagation, proper time, redshift, and local special relativity are exactly the corresponding GR/SR structures of that metric sector.

On the derivational side, the higher-derivative quartic branch has now cleared two earlier filters. First, its minimal quartic representative passes the flat-background \(k^0/k^2\) exact-closure screen while retaining a positive quartic scalar stiffness. Second, at the present action-level and flat-branch quadratic scope, the same branch preserves controlled matter coupling through the bridge metric alone. Those results narrow the next question considerably. The issue is no longer whether the branch has \emph{some} positive feature. The issue is whether it remains answerable to the exact relativistic benchmark already fixed elsewhere in the active sequence.

That is the task of the present note.

\paragraph{Framework and claim boundary.}
\TheoryProgram{} denotes the broader \Aether{} / \Aflow{} framework built on the claims that the \Aether{} is the underlying four-dimensional substrate of reality, the \Aflow{} is its intrinsic ordered motion, observed three-dimensional space is the local experiential slice of that deeper substrate, S-time is the experienced order of change arising from matter, light, and the \Aflow{}, and observed expansion is the three-dimensional appearance of deeper four-dimensional ordered motion. Within that framework, \TheoryName{} denotes the exact-closure theory in which the effective gravitational dynamics are adopted to be exactly Einsteinian with universal matter coupling. In the active sequence, \emph{adoption} means use of that established relativistic dynamics without claiming substrate derivation, while \emph{derivation} is reserved for a first-principles recovery from explicit substrate variables.

The present note stays within that boundary. It does not claim a completed derivation of GR from substrate variables. It asks a narrower question: does the minimal quartic higher-derivative branch preserve the benchmark weak-field, causal, clock, and local-SR structures of \TheoryName{} at the present disciplined scope, or does it leave an observer-accessible deformation already at that level?

\section{Recovery-Compatibility Gate}

\begin{definition}[Recovery-compatibility with the active exact-relativistic line]
At the present action-level and flat-branch quadratic scope, a candidate branch is said to be compatible with the active exact-relativistic recovery line if all four of the following hold.
\begin{enumerate}
    \item \textbf{Weak-field/Newtonian preservation through \(k^0/k^2\).} The reduced observer-accessible branch correction contributes no unsuppressed \(k^0\) or \(k^2\) scalar term in the weak-field bridge sector, so that the benchmark Poisson equation and Newtonian limit remain unchanged through those orders.
    \item \textbf{Single-metric null cone and proper-time rule.} Matter and light remain coupled only to one effective metric \(g_{\mu\nu}\), so that null propagation and timelike duration are still governed by
    \begin{equation}
    0=g_{\mu\nu}\,dx^\mu dx^\nu,
    \qquad
    d\tau^2=-\frac{1}{c^2}g_{\mu\nu}\,dx^\mu dx^\nu.
    \label{eq:nullproper}
    \end{equation}
    \item \textbf{Local SR preservation.} The local inertial limit already fixed in the relativistic-recovery manuscript remains the observer-accessible limit of the branch, with no second low-energy cone or preferred-frame matter metric.
    \item \textbf{Infrared suppression of the surviving quartic self-energy.} The only residual branch correction is metric-mediated and remains below observer-accessible relevance in a clearly stated infrared regime.
\end{enumerate}
\end{definition}

\begin{remark}
This definition is deliberately stronger than the earlier matter-coupling gate and deliberately narrower than a full curved-background theorem. The branch need not vanish identically. It must instead preserve the exact relativistic benchmark through the presently claimed orders and leave only a parametrically subleading residual metric correction in the intended observer-accessible regime.
\end{remark}

\section{Reduced Quartic Branch Data}

The present note uses only the minimal quartic branch data already established in the preceding manuscripts. The branch keeps the bridge metric
\begin{equation}
\BridgeMetric_{AB}=G_{AB}-2U_AU_B
\label{eq:bridge}
\end{equation}
and the universal matter route
\begin{equation}
S_{\mathrm{matter}}[\Xi^*\BridgeMetric,\psi]
=
S_{\mathrm{matter}}[g,\psi],
\qquad
g_{\mu\nu}=\Xi^*\BridgeMetric_{AB},
\label{eq:matterroute}
\end{equation}
while adjoining the quartic projected spatial operator
\begin{equation}
\Delta S_4
=
-\frac{\lambda_4}{2M_*^2}
\int_{\Sub} d^4X\,\sqrt{G}\,
\bigl(\Delta_\perp S\bigr)^2
\label{eq:S4}
\end{equation}
on the collapsed-scalar baseline
\begin{equation}
\kappa_S=\mu_{SI}=\mu_S^2=0,
\qquad
m_S^2>0,
\qquad
\widehat M_{IJ}\ \text{positive definite}.
\label{eq:baseline}
\end{equation}

Use the same admissible homogeneous bridge background as in the earlier quartic notes:
\begin{equation}
\bar G_{AB}=\delta_{AB},
\qquad
\bar U^A=\delta^A{}_0,
\qquad
\bar S=\sigma_0X^0,
\qquad
\bar Q^I=0,
\label{eq:background}
\end{equation}
for which
\begin{equation}
\bar\BridgeMetric_{AB}=\eta_{AB}=\mathrm{diag}(-1,1,1,1).
\label{eq:etabridge}
\end{equation}

\begin{proposition}[Only surviving observer-accessible branch correction]
On the decaying flat branch defined by \eqref{eq:baseline}, after eliminating the nonmetric auxiliary sectors \((q^I,\chi,V_T^i,\sigma)\), the reduced observer-accessible action has the form
\begin{equation}
S_{\mathrm{obs}}^{(2)}[h,\psi]
=
S_{\mathrm{EH}}^{(2)}[h]
+
S_{\mathrm{matter}}^{(2)}[h,\psi]
+
\Delta S_{\mathrm{metric},4}^{(2)}[h],
\label{eq:Sobs}
\end{equation}
with
\begin{equation}
\Delta\mathcal L_{\mathrm{metric},4}^{(2)}
=
-\frac{1}{2}
\Pi_\phi^{(4)}(\omega,\mathbf{k})\,|\phi(\omega,\mathbf{k})|^2
\label{eq:Deltametric}
\end{equation}
and
\begin{equation}
\Pi_\phi^{(4)}(\omega,\mathbf{k})
=
m_S^2\sigma_0^2
\frac{\mathcal B_4(k^2)}
{m_S^2\omega^2+\mathcal B_4(k^2)},
\qquad
\mathcal B_4(k^2)=\frac{\lambda_4}{M_*^2}k^4.
\label{eq:Piphi}
\end{equation}
No independent tensor or transverse-vector observer-accessible self-energy survives at this order.
\end{proposition}

\begin{proof}
The bridge curvature term remains the ordinary Einstein--Hilbert quadratic tensor sector because the branch leaves the bridge metric \eqref{eq:bridge} itself unchanged. The preceding quartic consistency note showed that the \(\chi\), \(V_T^i\), and \(q^I\) sectors remain auxiliary on the decaying flat branch and that the full nontrivial residual correction collapses onto the scalar channel \eqref{eq:Piphi}. The preceding matter-coupling note showed that this residual correction is metric-mediated, not a direct matter coupling to a second observer-accessible field. Together those two results give \eqref{eq:Sobs}--\eqref{eq:Piphi}.
\end{proof}

\begin{remark}
The present note therefore does not reopen the already settled matter-route question. Its burden is instead to determine whether the sole surviving metric-mediated scalar correction \eqref{eq:Piphi} is compatible with the exact relativistic recovery benchmark already fixed elsewhere in the active sequence.
\end{remark}

\section{Weak-Field and Newtonian Preservation Through \texorpdfstring{$k^0/k^2$}{k0/k2}}

The benchmark weak-field/Newtonian sector of \TheoryName{} is already fixed by the adopted Einsteinian metric theory:
\begin{equation}
\nabla^2\Phi=4\pi G\rho,
\label{eq:poisson}
\end{equation}
\begin{equation}
ds^2
=
-\left(1+2\frac{\Phi}{c^2}\right)c^2dt^2
+
\left(1-2\frac{\Phi}{c^2}\right)\delta_{ij}\,dx^i dx^j,
\label{eq:weakfield}
\end{equation}
and
\begin{equation}
\frac{d^2x^i}{dt^2}
=
-\partial_i\Phi
+ \mathcal O\!\left(\frac{v^2}{c^2},\frac{\Phi^2}{c^4}\right).
\label{eq:newtonian}
\end{equation}

The quartic branch must preserve this benchmark through the same \(k^0/k^2\) screen that already ruled out the earlier candidates.

\begin{proposition}[No \(k^0\) or \(k^2\) weak-field deformation]
For the minimal quartic branch,
\begin{equation}
\mathcal B_4(k^2)
=
\mathcal B_{4,0}
+
\mathcal B_{4,2}k^2
+
\mathcal B_{4,4}k^4,
\qquad
\mathcal B_{4,0}=0,
\qquad
\mathcal B_{4,2}=0,
\qquad
\mathcal B_{4,4}=\frac{\lambda_4}{M_*^2}.
\label{eq:B4coeffs}
\end{equation}
Therefore the reduced observer-accessible branch correction contributes no unsuppressed \(k^0\) or \(k^2\) deformation to the weak-field/Newtonian bridge sector.
\end{proposition}

\begin{proof}
Equation \eqref{eq:Piphi} shows that every nontrivial observer-accessible branch correction is controlled by the scalar operator \(\mathcal B_4(k^2)\). For the minimal quartic representative that operator is exactly \(\lambda_4k^4/M_*^2\). Its \(k^0\) and \(k^2\) coefficients therefore vanish identically, which is precisely the coefficient-level condition needed to preserve the benchmark weak-field/Newtonian sector through those orders.
\end{proof}

\begin{corollary}[GR weak-field/Newtonian benchmark survives the first recovery check]
At the present coefficient level, the quartic branch preserves the GR weak-field/Newtonian sector through the \(k^0/k^2\) benchmark screen. In particular, the branch leaves unchanged the Poisson equation \eqref{eq:poisson}, the standard weak-field line element \eqref{eq:weakfield}, and the Newtonian reduction \eqref{eq:newtonian} through those orders.
\end{corollary}

\begin{remark}
This is the precise meaning of the first positive recovery-compatibility result. The branch is not being claimed to vanish identically at all higher orders. It is being shown to leave the accepted weak-field/Newtonian sector untouched through the exact \(k^0/k^2\) filter that the active sequence uses as its first observer-level screen.
\end{remark}

\section{Single Metric Null Cone and Proper-Time Rule}

The next check is structural rather than coefficient-level. The active exact-relativistic line requires one and only one operational geometry. In \TheoryName{}, the benchmark causal and clock rules are
\begin{equation}
0=g_{\mu\nu}\,dx^\mu dx^\nu,
\label{eq:null}
\end{equation}
\begin{equation}
d\tau^2=-\frac{1}{c^2}g_{\mu\nu}\,dx^\mu dx^\nu.
\label{eq:propertime}
\end{equation}

\begin{proposition}[The quartic branch preserves the single operational metric]
At the present action-level and flat-branch quadratic scope, the quartic branch introduces no second observer-accessible matter metric, no independent low-energy null cone, and no alternative proper-time rule. Null propagation and timelike duration remain governed by \eqref{eq:null} and \eqref{eq:propertime} built from the same effective metric \(g_{\mu\nu}\).
\end{proposition}

\begin{proof}
By construction the branch keeps unchanged the observer map and matter route \eqref{eq:matterroute}. Matter and light therefore continue to couple to the same single pullback metric \(g_{\mu\nu}\). Proposition 1 shows that the only surviving branch correction after auxiliary elimination is the metric self-energy term \eqref{eq:Deltametric}; no separate observer-accessible tensor, vector, or scalar field is promoted to an independent matter argument. Accordingly, the operational causal and clock rules remain those of the single metric \(g_{\mu\nu}\) alone.
\end{proof}

\begin{corollary}[No second low-energy cone]
At the present scope, the quartic branch does not reintroduce the failure modes that would destroy exact closure: no second universal matter metric, no direct preferred-frame light cone, and no separate low-energy clock law.
\end{corollary}

\section{Preservation of the Local SR Limit}

The relativistic-recovery manuscript already fixed the local special-relativistic limit of \TheoryName{} through the local inertial conditions
\begin{equation}
g_{\mu\nu}(p)=\eta_{\mu\nu},
\qquad
\partial_\alpha g_{\mu\nu}(p)=0,
\label{eq:localinertial}
\end{equation}
which imply
\begin{equation}
ds^2
=
-c^2dt'^2+\delta_{ij}\,dx'^i dx'^j
+ \mathcal O((x'-p)^2).
\label{eq:localminkowski}
\end{equation}
In that local frame, timelike motion obeys
\begin{equation}
d\tau
=
dt'\sqrt{1-\frac{v^2}{c^2}}
=
\frac{dt'}{\Lor},
\qquad
\Lor\equiv \frac{1}{\sqrt{1-v^2/c^2}}.
\label{eq:localsr}
\end{equation}

\begin{proposition}[The quartic branch preserves the active local SR limit]
At the present scope, the quartic branch preserves the local Minkowski limit \eqref{eq:localminkowski}, local Lorentz symmetry of the tangent-space metric, and the ordinary local SR kinematics \eqref{eq:localsr}.
\end{proposition}

\begin{proof}
The local SR statement of the active recovery manuscript rests on two ingredients: one single observer-accessible metric and no second low-energy cone or preferred-frame matter metric. Proposition 4 supplies exactly those ingredients for the quartic branch at the present scope. Proposition 3 then ensures that the branch leaves the weak-field/Newtonian metric sector untouched through the \(k^0/k^2\) orders that feed the local inertial reduction. Therefore the operational local limit remains the locally Minkowskian metric sector already fixed in the active recovery line, with no extra observer-accessible structure available to spoil local Lorentz symmetry.
\end{proof}

\begin{remark}
This is the correct reading of the \Aether{} / \Aflow{} ontology at the present derivational stage. The quartic branch may modify the auxiliary substrate-side route to the effective metric, but at the observer level it does not insert a second cone, second clock law, or second local kinematics. In that sense it remains answerable to the exact local SR benchmark rather than deforming it.
\end{remark}

\section{Infrared Regime of Subleading Quartic Self-Energy}

The fourth check is the only one that requires a new regime statement rather than a direct structural inheritance argument. The remaining correction \eqref{eq:Piphi} is not zero. The question is therefore: in what infrared regime does that residual correction stay below observer-accessible relevance at the present scope?

\begin{definition}[Recovery-compatible infrared regime]
For the present note, the quartic branch is said to lie in its recovery-compatible infrared regime when both of the following hold:
\begin{equation}
k^2\ll M_*^2,
\label{eq:irhierarchy1}
\end{equation}
\begin{equation}
\varepsilon_4(\omega,\mathbf{k})
\equiv
\frac{\mathcal B_4(k^2)}{m_S^2\omega^2}
=
\frac{\lambda_4 k^4}{m_S^2M_*^2\omega^2}
\ll 1.
\label{eq:epsilon4}
\end{equation}
\end{definition}

\begin{remark}
Condition \eqref{eq:irhierarchy1} is the ordinary long-wavelength requirement expected for a higher-derivative completion. Condition \eqref{eq:epsilon4} is the branch-specific requirement that the surviving quartic scalar correction remain parametrically smaller than the existing time-derivative block. Equivalently,
\begin{equation}
\frac{\sqrt{\lambda_4}}{m_SM_*}\,k^2\ll |\omega|.
\label{eq:nonultrasoft}
\end{equation}
Thus the intended infrared regime is long-wavelength but not ultrasoft in the sense of sending \(\omega\) to zero faster than \(k^2\).
\end{remark}

\begin{proposition}[Infrared expansion of the quartic self-energy]
In the regime \eqref{eq:irhierarchy1}--\eqref{eq:epsilon4}, the surviving metric self-energy obeys
\begin{equation}
\Pi_\phi^{(4)}(\omega,\mathbf{k})
=
m_S^2\sigma_0^2\,\varepsilon_4(\omega,\mathbf{k})
+ \mathcal O\!\bigl(\varepsilon_4^2\bigr).
\label{eq:PiIR}
\end{equation}
Accordingly, the residual branch correction is parametrically subleading to the active exact-relativistic benchmark at the present scope.
\end{proposition}

\begin{proof}
Equation \eqref{eq:Piphi} may be rewritten as
\begin{equation}
\Pi_\phi^{(4)}
=
m_S^2\sigma_0^2
\frac{\varepsilon_4}{1+\varepsilon_4}.
\label{eq:PiEps}
\end{equation}
When \(\varepsilon_4\ll 1\), the geometric expansion of \((1+\varepsilon_4)^{-1}\) gives
\begin{equation}
\frac{\varepsilon_4}{1+\varepsilon_4}
=
\varepsilon_4+\mathcal O(\varepsilon_4^2),
\end{equation}
which yields \eqref{eq:PiIR}. Since the preceding sections already showed that the branch does not alter the \(k^0/k^2\) weak-field benchmark, the single null cone/proper-time rule, or the local SR limit, a residual metric correction suppressed by \(\varepsilon_4\) is below observer-accessible relevance in the present recovery-compatibility sense.
\end{proof}

\begin{remark}
The suppression statement is not uniform in the ultrasoft static corner \(\omega\to 0\) taken before the quartic hierarchy \eqref{eq:epsilon4} is enforced. The present note therefore does not claim that every static or arbitrarily low-frequency configuration has already been shown harmless. It claims only the cleaner and presently sufficient statement: in the long-wavelength, non-ultrasoft infrared regime \eqref{eq:irhierarchy1}--\eqref{eq:epsilon4}, the surviving quartic correction is parametrically subleading and does not upset the active recovery benchmark.
\end{remark}

\section{Recovery-Compatibility Result}

\begin{theorem}[Quartic branch recovery-compatibility gate]
Assume the minimal quartic branch \eqref{eq:bridge}--\eqref{eq:baseline} on the decaying flat branch, together with the recovery-compatible infrared regime \eqref{eq:irhierarchy1}--\eqref{eq:epsilon4}. Then, at the present action-level and flat-branch quadratic scope, the quartic branch is compatible with the active exact-relativistic recovery line.
\end{theorem}

\begin{proof}
Definition 1 requires four checks. The first follows from Proposition 2 and Corollary 1: the GR weak-field/Newtonian sector is preserved through the \(k^0/k^2\) benchmark screen. The second follows from Proposition 3 and Corollary 2: the branch preserves the single operational metric, the unique null cone, and the proper-time rule. The third follows from Proposition 4: the active local SR limit remains the observer-accessible local limit of the branch. The fourth follows from Definition 2 and Proposition 5: in the intended infrared regime, the only surviving residual correction is a metric-mediated quartic self-energy suppressed by \(\varepsilon_4\). All four conditions of Definition 1 therefore hold.
\end{proof}

\begin{corollary}[Current promotion rule]
The quartic branch may be carried forward to a later local curved-background continuation only in the same disciplined sense established here: the branch is promoted as recovery-compatible on the long-wavelength, non-ultrasoft infrared regime, not as a fully unconstrained global deformation and not as a completed derivation.
\end{corollary}

\section{What This Note Establishes and What It Does Not}

The present note establishes the following.
\begin{enumerate}
    \item It states a dedicated recovery-compatibility gate for the quartic branch in exactly the order now required by the active manuscript program.
    \item It shows that the quartic branch preserves the GR weak-field/Newtonian sector through the benchmark \(k^0/k^2\) screen.
    \item It shows that the branch preserves the single metric null cone and proper-time rule already fixed in the dynamics manuscript.
    \item It shows that the local SR limit fixed in the relativistic-recovery manuscript remains the observer-accessible local limit of the branch.
    \item It states the infrared regime in which the only surviving quartic metric self-energy is parametrically subleading and below observer-accessible relevance in the present exact-closure sense.
    \item It records the corresponding current-scope positive result: the quartic branch is compatible with the active exact-relativistic recovery line in that regime.
\end{enumerate}

The present note does \emph{not} establish the following.
\begin{enumerate}
    \item It does not prove a full nonlinear Einsteinian derivation from substrate variables.
    \item It does not claim uniform suppression of the quartic self-energy in the ultrasoft static corner.
    \item It does not yet carry out the corresponding local curved-background continuation.
    \item It does not settle nonlinear stability or any global admissible-background question.
\end{enumerate}

\begin{remark}
These limits matter scientifically. The note is positive, but its positivity is disciplined. It clears the exact next gate needed by the current branch sequence and no more.
\end{remark}

\section{Conclusion}

The higher-derivative quartic branch was not allowed to move directly from a positive flat-background coefficient test and a positive matter-coupling gate to a curved-background promotion. It first had to be checked against the active exact-relativistic recovery line. The present note now supplies that check at the current disciplined scope.

The answer is positive in exactly the sense the active sequence now needs. The branch preserves the GR weak-field/Newtonian sector through the \(k^0/k^2\) screen, preserves the single operational metric together with its null cone and proper-time rule, preserves the local SR limit already fixed in the relativistic-recovery manuscript, and leaves only a quartic metric self-energy that is parametrically subleading in the long-wavelength, non-ultrasoft infrared regime.

That is enough to justify the next local continuation step, but only with the same claim boundary carried forward explicitly. The quartic branch is not yet promoted as a completed derivation and not yet promoted beyond the regime in which its residual self-energy is demonstrably subleading. The next note, if undertaken, should therefore be a local curved-background continuation that keeps the present recovery-compatible hierarchy visible rather than silently enlarging the claim.

\end{document}
